\subsection{SUBMODULES; QUOTIENT MODULES}

Let E be an A-module and M a subset of E; for the A-module structure on E to induce an A-module structure on M, it is necessary and sufficient that M be a stable subgroup of E (I, § 4, no. 3), for if this is so, the structure of a group with operators induced on M obviously satisfies axioms (M \(_{II}\) ), (M \(_{III}\) ) and (M \(_{IV}\) ); then M with this structure (or, by an abuse of language, the set M itself) is called a submodule of E; the canonical injection M → E is a linear mapping. When E is a vector space, its submodules are called vector subspaces (or simply subspaces if no confusion can arise).

Examples. (1) In any module E the set consisting of 0 is a submodule (the zero submodule, often denoted by 0, by an abuse of notation).

\begin{enumerate}
\def\labelenumi{(\arabic{enumi})}
\setcounter{enumi}{1}
\item
  Let \(\mathbf{A}\) be a ring. The submodules of \(\mathbf{A}_s\) (resp. \(\mathbf{A}_d\) ) are just the left ideals (resp. right ideals) of the ring \(\mathbf{A}\) .
\item
  Let E be an A-module, x an element of E and a a left ideal of A. The set of elements \(\alpha x\) , where \(\alpha\) runs through a, is a submodule of E, denoted by ax.
\item
  In a commutative group G considered as a Z-module (no. 1), every subgroup of G is also a submodule.
\end{enumerate}

*(5) Let I be an open interval of the real line \(\mathbf{R}\) ; the set C of real-valued functions defined and continuous on I is a vector subspace of the vector space \(\mathbf{R}^{\mathrm{I}}\) of all real-valued functions defined on I. Similarly, the set D of differentiable functions on I is a vector subspace of \(\mathbf{C}\).*

Let E be an A-module. Every equivalence relation compatible (I, §1, no. 6) with the module structure on E is of the form \(x - y \in \mathbf{M}\) , where M is a stable subgroup of E (I, § 4, no. 4), that is a submodule of E. It is immediately verified that the structure of a group with operators on the quotient group E/M (I, § 4, no. 4) is an A-module structure, under which the canonical mapping E → E/M is linear; with this structure, E/M is called the quotient module of E by the submodule M. A quotient module of a vector space E is called a quotient vector space (or simply quotient space) of E.

Example (6). Every left ideal a in a ring A defines a quotient module \(A_{s}/a\) of the left A-module \(A_{s}\) ; by an abuse of notation, this quotient module is often denoted by A/a.

Let E, F be two A-modules. It follows from the general properties of groups with operators (I, § 4, no. 5) (or directly from the definitions) that if \(u: \mathbf{E} \to \mathbf{F}\) is a linear mapping, the image under \(u\) of any submodule of E is a submodule of F and the inverse image under \(u\) of any submodule of F is a submodule of

E. In particular, the kernel \(N = u^{-1}(0)\) is a submodule of E and the image \(u(\mathbf{E})\) of E under u is a submodule of F (I, §4, no. 6, Proposition 7); by an abuse of language \(u(\mathbf{E})\) is called the image of u. The quotient module E/N is also called the coimage of u and the quotient module \(F/u(\mathbf{E})\) the cokernel of u. In the canonical decomposition of u (I, §4, no. 5)

\[
u: \mathrm{E} \xrightarrow {p} \mathrm{E} / \mathrm{N} \xrightarrow {v} u (\mathrm{E}) \xrightarrow {i} \mathrm{F}\tag{11}
\]

v is an isomorphism of the coimage of u onto the image of u (no. 2, Proposition 2). For u to be injective, it is necessary and sufficient that its kernel be zero; for u to be surjective, it is necessary and sufficient that its cokernel be zero.

The kernel, image, coimage and cokernel of u are denoted respectively by Ker u, Im u, Coim u, Coker u.

Remark. Let M be a submodule of an A-module E and \(\phi: \mathrm{E} \to \mathrm{E} / \mathrm{M}\) the canonical homomorphism. For an A-linear mapping \(u: \mathrm{E} \to \mathrm{F}\) to be of the form \(v \circ \phi\) , where \(v\) is a linear mapping of \(\mathrm{E} / \mathrm{M}\) into \(\mathbf{F}\) , it is necessary and sufficient that \(\mathrm{M} \subset \operatorname{Ker}(u)\) ; for, if this condition holds, the relation \(x - y \in \mathbf{M}\) implies \(u(x) = u(y)\) , hence \(u\) is compatible with this equivalence relation and clearly the mapping \(v: \mathrm{E} / \mathrm{M} \to \mathrm{F}\) derived from \(u\) by taking quotients is linear.

